\subsection*{Historical cluster-aware preference study}
\label{subsec:preference}

The 01549 first-round revision reported a blinded three-alternative forced-choice study on 30 objects with 24 participants. Participants compared conservative fixed-scale, aggressive fixed-scale, and C3 outputs for texture naturalness, shape consistency, and overall quality; condition order was randomized per trial. Each criterion had 720 first-choice decisions. Table~\ref{tab:preference} reports the aggregate counts; participant-level responses are not redistributed with this revision.

\begin{table}[t]
\centering\scriptsize
\caption{First-round 01549 3AFC results retained as reported. Percentages use 720 decisions per criterion; preference is not a measure of reference texture fidelity.}
\label{tab:preference}
\setlength{\tabcolsep}{3pt}
\begin{tabular}{lccc}
\toprule
Method & Texture naturalness & Shape consistency & Overall quality \\
\midrule
$s=1.25$ & 332/720 (46.1\%) & 96/720 (13.3\%) & 176/720 (24.4\%) \\
$s=2.50$ & 104/720 (14.4\%) & 331/720 (46.0\%) & 126/720 (17.5\%) \\
C3 (TCAS) & 284/720 (39.4\%) & 293/720 (40.7\%) & 418/720 (58.1\%) \\
\bottomrule
\end{tabular}
\end{table}

The reported analysis used a two-level participant/object bootstrap with 10,000 resamples. C3's overall-quality share was 58.1\% (95\% CI [50.1\%, 65.8\%]); its differences from the conservative and aggressive baselines were +33.6 points (CI [19.3, 47.2]) and +40.6 points (CI [29.4, 51.8]). The single-criterion advantages for texture naturalness and shape consistency had intervals crossing zero. This bounded preference result does not validate texture fidelity, unseen-view quality, or the Fresh C cohort.

